% Source: https://artofproblemsolving.com/community/c3h2793863_graph_of_equation
% Stable identity: aops-topic-2793863-q01
$(x+3)y=(x-3)(x+3)$ means $(x+3)(y-x+3)=0$: the graph is the union of the line $y=x-3$ and the vertical line $x=-3$. (Dividing by $x+3$ loses the second line.)
\begin{asy}
usepackage("amsmath");
size(4cm);
defaultpen(fontsize(10pt));
draw((-5,0)--(5,0),Arrow);draw((0,-7)--(0,3),Arrow);draw((-4,-7)--(5,2));draw((-3,-7)--(-3,3));dot((0,-3));dot((3,0));dot((-3,0));label("$(0,-3)$",(0,-3),E);label("$(3,0)$",(3,0),SE);label("$(-3,0)$",(-3,0),NW);
\end{asy}


$(0,a)$ is on $y=x-3$: $a=-3$. The points on the $x$-axis are $(3,0)$ and $(-3,0)$. So $a+b+c=-3+3-3=\boxed{-3}$.
